% =====================================================================
% SECCIÓN 4: MODELADO Y ARQUITECTURA
% =====================================================================

% ---------------------------------------------------------------------
% Diapositiva 10: Jerarquía de Modelos Supervisados
% ---------------------------------------------------------------------
\section{Modelado}

\begin{frame}{Jerarquía de modelos supervisados}
  \framesubtitle{Tres niveles de complejidad representacional y diseño experimental $2 \times 2$}

  \begin{columns}[T]
    \column{0.32\textwidth}
    \ugrblock{1. Líneas base léxicas}{
      {\scriptsize
      \begin{itemize}
        \item \textbf{TF-IDF:} Unigramas y bigramas ($V = 15\,000$).
        \item \textbf{Regresión Logística:} Regularización $L_2$ ($C=1{,}0$).
        \item \textbf{SVM lineal:} Margen con norma $L_2$.
      \end{itemize}}
    }

    \column{0.33\textwidth}
    \ugrblock{2. Modelo combinado}{
      {\scriptsize
      \begin{itemize}
        \item \textbf{Enriquecimiento:} TF-IDF + 19 rasgos psicolingüísticos.
        \item \textbf{Clasificador:} LR calibrada.
        \item \textbf{Objetivo:} Evaluar si rasgos expertos superan n-gramas.
      \end{itemize}}
    }

    \column{0.33\textwidth}
    \ugrblockok{3. Modelo contextual}{
      {\scriptsize
      \begin{itemize}
        \item \textbf{BETO:} Transformer preentrenado (\texttt{dccuchile}).
        \item \textbf{Capacidad:} 110M params (\emph{fine-tuning}).
        \item \textbf{Atención:} Orden, polisemia y negaciones sintácticas.
      \end{itemize}}
    }
  \end{columns}

  \vskip-0.10cm
  \ugrblock{Diseño experimental $2 \times 2$ (Mono vs.\ Multi)}{
    {\scriptsize Cada familia se entrenó en dos modalidades: \textbf{Mono} (solo Twitter nativo) y \textbf{Multi} (Twitter + Reddit traducido), evaluadas sobre el conjunto nativo de prueba común.}
  }

  \note{
    Planteamos una progresión metodológica rigurosa: líneas base léxicas puras (TF-IDF con Regresión Logística y SVM); un modelo combinado que incorpora los 19 rasgos psicolingüísticos; y el modelo contextual de vanguardia BETO. Cada modelo se entrena en modalidad mono y multi para comparar el impacto del aumento de datos.
  }
\end{frame}

% ---------------------------------------------------------------------
% Diapositiva 11: Arquitectura de BETO y Configuración Experimental
% ---------------------------------------------------------------------
\begin{frame}{Arquitectura de BETO y configuración experimental}
  \framesubtitle{Ajuste fino supervisado con restricciones computacionales reales}

  \begin{columns}[T]
    \column{0.48\textwidth}
    \begin{center}
      \resizebox{0.88\columnwidth}{!}{% Arquitectura del ajuste fino de BETO para clasificación binaria.
\begin{tikzpicture}[font=\footnotesize, node distance=0.55cm]
    \node[proceso=2.6cm, fill=ugrGris!8, draw=ugrGris] (tx) {\textbf{Texto}\\ \emph{``no me}\\ \emph{siento bien''}};
    \node[proceso=2.6cm, right=of tx] (tok) {\textbf{Tokenización}\\ WordPiece\\ \texttt{[CLS]} \dots{} \texttt{[SEP]}\\ máx.\ 128 tokens};
    \node[proceso=2.9cm, right=of tok, minimum height=2.6cm, fill=ugrAzul!10, line width=0.9pt] (enc)
        {\textbf{BETO}\\ codificador Transformer\\ 12 capas \texttimes{} 12 cabezas\\ $d = 768$\\ {\color{ugrGris}\scriptsize preentrenado, se ajusta}};
    \node[proceso=2.2cm, right=of enc] (cls) {\textbf{Vector}\\ $h_{\texttt{[CLS]}} \in \mathbb{R}^{768}$};
    \node[proceso=2.2cm, right=of cls, draw=ugrNaranja, fill=ugrNaranja!8] (lin) {\textbf{Capa lineal}\\ $768 \to 2$\\ + softmax};
    \node[resultado=2.9cm, below=0.7cm of lin] (out) {$P(\text{depresivo}\mid x)$};

    \draw[flecha] (tx) -- (tok);
    \draw[flecha] (tok) -- (enc);
    \draw[flecha] (enc) -- (cls);
    \draw[flecha] (cls) -- (lin);
    \draw[flecha] (lin) -- (out);
    \node[below=0.15cm of enc, font=\scriptsize, text=ugrGris, align=center] {AdamW, $\eta = 2\times10^{-5}$, lote 16, FP16};
\end{tikzpicture}
}
      \vskip0.05cm
      {\tiny\color{ugrGris}Diagrama: WordPiece $\to$ BETO ($d=768$) $\to$ Capa lineal $\to$ Softmax.\par}
    \end{center}

    \column{0.50\textwidth}
    \ugrblock{Configuración de hiperparámetros y hardware}{
      {\scriptsize
      \begin{itemize}
        \item \textbf{Optimizador:} AdamW ($\text{weight decay} = 0{,}01$, $\eta = 2 \times 10^{-5}$).
        \item \textbf{Secuencia y lote:} Longitud máx.\ 128 tokens, lote de 16.
        \item \textbf{Entrenamiento:} 3 épocas con \emph{early stopping} (óptimo en época 2).
        \item \textbf{Hardware:} GPU NVIDIA RTX 3060 Ti (8\,GB VRAM) en FP16.
      \end{itemize}}
    }
  \end{columns}

  \note{
    BETO procesa la secuencia tokenizada con WordPiece; el vector contextual de la salida CLS en la capa 12 alimenta una capa lineal hacia dos neuronas con softmax. El ajuste fino se completó en hardware local mediante precisión mixta FP16, alcanzando el óptimo en la época 2.
  }
\end{frame}

% ---------------------------------------------------------------------
% Diapositiva 12: Ejemplo de Clasificación End-to-End
% ---------------------------------------------------------------------
\begin{frame}{Ejemplo de Clasificación End-to-End}
  \framesubtitle{Flujo secuencial de inferencia sobre una publicación individual}

  \vskip-0.1cm
  \begin{center}
  \resizebox{0.96\textwidth}{!}{
  \begin{tikzpicture}[
    paso/.style={rectangle, rounded corners=4pt, draw=ugrAzul, fill=ugrFondoClaro, line width=0.8pt, inner sep=6pt, align=center, text width=3.0cm, font=\small},
    pasoOk/.style={rectangle, rounded corners=4pt, draw=ugrVerde, fill=ugrVerde!10, line width=0.8pt, inner sep=6pt, align=center, text width=3.0cm, font=\small},
    flechaFlujo/.style={-{Stealth[length=2.5mm, width=1.8mm]}, line width=1.2pt, color=ugrCeleste}
  ]
    % Nodos
    \node[paso] (n1) at (0, 0) {
      \textbf{\color{ugrAzul}1. Texto de entrada}\\[3pt]
      {\scriptsize Publicación individual}\\[4pt]
      \emph{\tiny ``Nunca me sentí tan vacío, no le encuentro sentido a nada...''}
    };

    \node[paso] (n2) at (3.8, 0) {
      \textbf{\color{ugrAzul}2. Tokenización}\\[3pt]
      {\scriptsize WordPiece ($V = 31\text{k}$)}\\[4pt]
      \texttt{\tiny [CLS] nunca me sentí tan vacío , no le encuentro ... [SEP]}\\[3pt]
      {\tiny Longitud $L = 128$}
    };

    \node[paso] (n3) at (7.6, 0) {
      \textbf{\color{ugrAzul}3. Modelo BETO}\\[3pt]
      {\scriptsize Transformer Bidireccional}\\[4pt]
      {\tiny 12 capas $\cdot$ 110M params}\\[2pt]
      {\tiny Vector contextual}\\[2pt]
      \textbf{\tiny $\mathbf{h}_{\text{[CLS]}} \in \mathbb{R}^{768}$}
    };

    \node[pasoOk] (n4) at (11.4, 0) {
      \textbf{\color{ugrVerde}4. Predicción}\\[3pt]
      {\scriptsize Capa lineal + Softmax}\\[4pt]
      {\tiny $P(\text{Depresión}) = \mathbf{0{,}89}$}\\[2pt]
      {\tiny $P(\text{Control}) = 0{,}11$}\\[3pt]
      \textbf{\footnotesize\color{ugrRojo}Indicios de depresión}
    };

    % Flechas
    \draw[flechaFlujo] (n1) -- (n2);
    \draw[flechaFlujo] (n2) -- (n3);
    \draw[flechaFlujo] (n3) -- (n4);
  \end{tikzpicture}
  }
  \end{center}

  \vskip0.05cm
  \begin{columns}[T]
    \column{0.49\textwidth}
    \ugrblock{Captura contextual vs.\ bolsa de palabras}{
      BETO preserva el orden secuencial y resuelve la negación (\emph{``no le encuentro sentido''}), que los modelos léxicos dispersan en términos inconexos.
    }
    \column{0.49\textwidth}
    \ugrblockok{Salida probabilística calibrada}{
      La distribución de salida provee una probabilidad posterior interpretable, permitiendo ajustar umbrales de decisión según la sensibilidad clínica requerida.
    }
  \end{columns}

  \note{
    Ilustramos el pipeline de clasificación end-to-end: desde el microtexto crudo hasta la probabilidad final. El texto pasa por la tokenización WordPiece con tokens especiales CLS y SEP; las 12 capas de auto-atención de BETO generan el embedding contextual de 768 dimensiones, y la capa lineal clasifica la publicación con una probabilidad de indicios de depresión del 89\%.
  }
\end{frame}
